% Part: counterfactuals
% Chapter: introduction
% Section: counterfactuals

\documentclass[../../../include/open-logic-section]{subfiles}

\begin{document}

\olfileid[tr]{cnt}{int}{cnt}

\olsection{Karşıolgusallar}

İngilizcedeki ``if \dots{} then \dots'' yapılarının çok yaygın ve önemli bir
biçimi, \emph{to be} fiilinin geçmiş dilek kipiyle kurulur: ``Eğer durum
\dots{} olsaydı \dots{} olurdu.'' Böyle bir koşulun önbileşeni genellikle
yanlış, yani olguya aykırı olduğundan bunlara \emph{karşıolgusal koşullar}
(\emph{to be} fiilinin dilek kipini kullandıkları için ayrıca \emph{dilek
kipli koşullar}) denir. Bunlar ``Eğer durum \dots{} ise durum \dots{}tir''
biçimindeki \emph{bildirme kipli} koşullardan ayrılır. Karşıolgusal ve
bildirme kipli koşulların doğruluk koşulları farklıdır. Adams'ın ünlü
örneğini düşünün:
\begin{quote}
  Oswald Kennedy'yi öldürmediyse başka biri öldürdü.

  Oswald Kennedy'yi öldürmemiş olsaydı başka biri öldürürdü.
\end{quote}
Birincisi bildirme kipli, ikincisi karşıolgusaldır. Birincisi açıkça doğrudur:
Başkan John F. Kennedy'nin \emph{biri} tarafından öldürüldüğünü biliyoruz;
eğer bu kişi (Warren Raporu'nun tersine) Lee Harvey Oswald değilse Kennedy'yi
başka biri öldürmüştür. İkincisi başka bir şey söyler. Oswald Kennedy'yi
öldürmemiş olsaydı, yani Dallas'taki silahlı saldırı önlenmiş ya da başarısız
olmuş olsaydı, tarihin daha sonra başka bir suikastın başarılı olacağı biçimde
ilerleyeceğini öne sürer. Bunun doğru olması için güçlü odakların JFK'nin
ölümünü güvence altına almak üzere komplo kurmuş olması gerekirdi (birçok JFK
komplo kuramcısının inandığı gibi).

\emph{Bildirme kipli} koşulun maddi koşulla doğru biçimde yakalanıp
yakalanmadığı---özellikle maddi koşulun paradokslarının, İngilizcedeki
bildirme kipli koşulların doğruluk koşullarını yine de vermesiyle bağdaşacak
biçimde ``açıklanıp'' açıklanamayacağı---hâlâ tartışmalıdır. Buna karşılık
karşıolgusal koşulların maddi koşullarla doğru biçimde simgeleştirilemeyeceği
tartışmasızdır. Bu açıktır; çünkü karşıolgusalların önbileşenleri genellikle
yanlış olsa da önbileşeni yanlış olan her karşıolgusal doğru değildir. Örneğin
Warren Raporu'na inanıyor ve JFK'ye suikast düzenlemek için komplo
kurulmadığını düşünüyorsanız Adams'ın karşıolgusal koşulu buna örnektir.

Karşıolgusal koşullar nedensel akıl yürütmede önemli bir rol oynar:
karşıolgusalların başlıca kullanımlarından biri nedensel bağıntıları ifade
etmektir. Örneğin bir kibriti çakmak onun yanmasına neden olur; bunu ``Bu
kibrit çakılsaydı yanardı'' diyerek ifade edebilirsiniz. Maddi koşullar ve
genel olarak bildirme kipli koşullar bunu ifade edemez: ``kibrit çakılır
$\lif$ kibrit yanar'', kibrit hiç çakılmazsa, çakılsaydı ne olacağından
bağımsız olarak doğrudur. Daha kötüsü, ``kibrit çakılır $\lif$ kibrit bir
çiçek demetine dönüşür'' de kibrit hiç çakılmazsa doğrudur; oysa kibrit
çakılsaydı kesinlikle bir çiçek demetine dönüşmezdi.

Karşıolgusalların doğru mantığının tam olarak ne olduğu hâlâ tartışılmaktadır.
Stalnaker ile Lewis karşıolgusalların etkili bir çözümlemesini sundu. Onlara
göre ``$S$ olsaydı $T$ olurdu'' karşıolgusalı, ancak ve ancak gerçek dünyanın
durumuna en yakın ve $S$'nin doğru olduğu karşıolgusal durumda (``olası
dünyada'') $T$ doğruysa doğrudur. Olası dünyalar ontolojisine gönderimde
bulunduğundan buna ``ontik'' çözümleme denir. Başka çözümlemeler koşullu
olasılıkları ya da inanç revizyonu kuramlarını kullanır. Karşıolgusallar için
önerilen pek çok farklı mantık vardır. Tek bir Lewis--Stalnaker karşıolgusal
mantığı bile yoktur: Stalnaker ile Lewis en yakın olası dünyalara gönderim
yapan benzer açıklamalar önermiş olsalar da benimsedikleri varsayımlar farklı
geçerli çıkarımlara yol açar.

\end{document}
